\documentclass[12pt]{article}
\usepackage[utf8]{inputenc}
\usepackage[T1]{fontenc}
\usepackage[margin=1in]{geometry}
\usepackage{setspace}
\usepackage{xcolor}
\usepackage[colorlinks=true, urlcolor=blue!60!black, linkcolor=blue!60!black]{hyperref}
\setlength{\parindent}{2em}
\setlength{\parskip}{1.25em}
\renewcommand{\baselinestretch}{1.0}
\begin{document}
\begin{center}
{\LARGE \textit{Murphy v. National Collegiate Athletic Association}, 584 U.S. 453 (2018)\par}
\vspace{0.8em}
{\large Case Brief\par}
\end{center}
\vspace{1em}

\section{Memory Jogger}

In \href{https://www.courtlistener.com/opinion/4497655/murphy-v-national-collegiate-athletic-assn/}{\textit{Murphy v. National Collegiate Athletic Association}}, the Supreme Court held that the Tenth Amendment's anti-commandeering rule bars Congress from issuing direct orders to state legislatures---whether affirmative commands to enact laws or negative prohibitions forbidding states from repealing or modifying existing laws---and struck down the federal Professional and Amateur Sports Protection Act (PASPA) in its entirety.

\section{Detailed Case Facts}

In 1992, Congress enacted the Professional and Amateur Sports Protection Act (PASPA), 28 U.S.C. \S\S~3701--3704, to stem the spread of state-sanctioned sports gambling and safeguard the integrity of amateur and professional sports (\href{https://supreme.justia.com/cases/federal/us/584/16-476/}{Justia Opinion}). PASPA's central provision, 28 U.S.C. \S~3702(1), made it unlawful for a state or any of its political subdivisions to sponsor, operate, advertise, promote, license, or authorize by law or compact any lottery, sweepstakes, or other betting, gambling, or wagering scheme based on competitive sporting events. A companion provision, \S~3702(2), barred private parties from sponsoring, operating, advertising, or promoting such schemes if done pursuant to the law or compact of a governmental entity. Crucially, PASPA did not criminalize sports gambling under federal law; instead, \S~3703 authorized the U.S. Attorney General and amateur or professional sports organizations to bring civil actions for injunctive relief (\href{https://www.law.cornell.edu/supremecourt/text/16-476}{Cornell LII}).

PASPA included grandfather clauses in \S~3704(a)(1)--(2) permitting existing sports-gambling operations to continue in four states---Nevada, Delaware, Oregon, and Montana---and a special one-year window in \S~3704(a)(3) allowing New Jersey to legalize sports betting in Atlantic City casinos within one year of PASPA's effective date (\href{https://www.law.cornell.edu/supremecourt/text/16-476}{Cornell LII}). New Jersey did not act during that window. Two decades later, facing casino competition from neighboring states and tribal lands, New Jersey voters in 2011 approved a state constitutional amendment empowering the legislature to legalize sports wagering at Atlantic City casinos and horseracing tracks. The legislature enacted a 2012 statute affirmatively legalizing and regulating sports betting (the 2012 Act). After federal courts enjoined the 2012 Act, the legislature passed a 2014 statute (the 2014 Act) framed as a partial repealer: it repealed state prohibitions on sports wagering for persons aged 21 and older at Atlantic City casinos and New Jersey horseracing tracks, while leaving bans intact for wagers involving New Jersey college teams or collegiate events in the state (\href{https://supreme.justia.com/cases/federal/us/584/16-476/}{Justia Opinion}).

\section{Procedural History}

After New Jersey passed the 2012 Act, the National Collegiate Athletic Association (NCAA) and four major professional sports leagues (the NBA, NFL, NHL, and MLB) sued Governor Chris Christie and other state officials in the U.S. District Court for the District of New Jersey (\textit{Christie I}). The District Court rejected New Jersey's anti-commandeering defense and enjoined the 2012 Act, 926 F. Supp. 2d 551 (D.N.J. 2013). A divided Third Circuit panel affirmed, 730 F.3d 208 (3d Cir. 2013), concluding that PASPA imposed no affirmative mandate and did not bar states from repealing existing prohibitions (\href{https://supreme.justia.com/cases/federal/us/584/16-476/}{Justia Opinion}). The Supreme Court denied certiorari in June 2014, 573 U.S. 931 (2014).

When New Jersey enacted the 2014 partial-repeal Act, the leagues sued again (\textit{Christie II}), and the New Jersey Thoroughbred Horsemen's Association (NJTHA), operator of Monmouth Park Racetrack, joined the defense (\href{https://www.scotusblog.com/2018/05/opinion-analysis-justices-strike-down-federal-sports-gambling-law/}{SCOTUSblog}). The District Court again enjoined the state statute, 61 F. Supp. 3d 488 (D.N.J. 2014). A Third Circuit panel affirmed, 799 F.3d 259 (3d Cir. 2015), and the en banc Third Circuit affirmed 9--3, 832 F.3d 389 (3d Cir. 2016), holding that selectively repealing bans to channel wagering to specific venues amounted to ``authorization'' under PASPA and did not violate the Tenth Amendment because PASPA required no affirmative state action (\href{https://www.law.cornell.edu/supremecourt/text/16-476}{Cornell LII}). The Supreme Court granted certiorari on June 27, 2017, consolidating the petitions of New Jersey (No. 16-476) and NJTHA (No. 16-477). Following Phil Murphy's election as Governor, the case was recaptioned \textit{Murphy v. National Collegiate Athletic Association}. The Court heard argument on December 4, 2017, and decided the case on May 14, 2018 (\href{https://supreme.justia.com/cases/federal/us/584/16-476/}{Justia Opinion}).

\section{Judicial Votes}

The Supreme Court reversed the Third Circuit, splitting \textbf{7--2} on the Tenth Amendment anti-commandeering holding and \textbf{6--3} on severability (\href{https://supreme.justia.com/cases/federal/us/584/16-476/}{Justia Opinion}; \href{https://reason.com/volokh/2018/05/14/sports-gambling-decision-is-a-major-vict/}{Volokh Conspiracy}):

\begin{itemize}
    \item \textbf{Majority Opinion (Parts I--V and VI--A, 7--2; Part VI--B, 6--3):} Justice Samuel Alito delivered the opinion of the Court, joined in full by Chief Justice John Roberts and Justices Anthony Kennedy, Clarence Thomas, Elena Kagan, and Neil Gorsuch. Justice Stephen Breyer joined all parts except Part VI--B (severability).
    \item \textbf{Concurring Opinion:} Justice Clarence Thomas joined the majority opinion in full and filed a separate concurrence questioning modern severability doctrine.
    \item \textbf{Concurring in Part and Dissenting in Part:} Justice Stephen Breyer joined Parts I--V and VI--A (agreeing that \S~3702(1)'s anti-authorization and anti-licensing provisions violate the anti-commandeering rule) but dissented from Part VI--B, joining Justice Ginsburg's dissent in part on the severability of \S~3702(2).
    \item \textbf{Dissenting Opinion:} Justice Ruth Bader Ginsburg filed a dissenting opinion, joined in full by Justice Sonia Sotomayor and joined in part by Justice Stephen Breyer.
\end{itemize}

\section{Holding}

\textbf{Judgment:} Reversed. PASPA was invalidated in its entirety (\href{https://www.law.cornell.edu/supremecourt/text/16-476}{Cornell LII}).

\textbf{Rule Announced:} Under the Constitution's system of dual sovereignty reflected in the Tenth Amendment, Congress lacks authority to issue direct orders to state legislatures, whether by affirmatively compelling them to enact or enforce a federal regulatory program or by prohibiting them from enacting, modifying, or repealing state laws. Moreover, a valid preemption provision under the Supremacy Clause must both exercise an enumerated congressional power and be best read as regulating private actors (by conferring federal rights or imposing federal restrictions on individuals), rather than commanding state governments. Because PASPA's bars on state ``authorization'' and ``licensing'' in 28 U.S.C. \S~3702(1) issue direct orders to state legislatures without regulating private actors, they violate the anti-commandeering rule; and because no remaining provision of PASPA is severable, the entire statute is invalid (\href{https://supreme.justia.com/cases/federal/us/584/16-476/}{Justia Opinion}).

\textbf{Holding vs. Dicta:} The binding holdings are (1) that a state's total or partial repeal of laws banning sports gambling constitutes ``authoriz[ation]'' under PASPA \S~3702(1); (2) that \S~3702(1)'s prohibitions on state authorization and licensing violate the Tenth Amendment's anti-commandeering rule regardless of whether framed as affirmative mandates or negative prohibitions; (3) that valid federal preemption operates only through laws regulating private actors; and (4) that PASPA's remaining provisions---including \S~3702(2)'s restriction on private operations ``pursuant to'' state law and advertising bans---are nonseverable. By contrast, Justice Alito's remark that Congress may regulate sports gambling directly under the Commerce Clause was unnecessary to the judgment and therefore constitutes dicta---a point highlighted by Justice Thomas's concurrence expressing doubt about federal power over purely intrastate wagering (\href{https://www.law.cornell.edu/supremecourt/text/16-476}{Cornell LII}).

\section{Analysis of Opinions}

\subsection{Majority Opinion (Justice Alito)}

Justice Alito first resolved the statutory meaning of \S~3702(1), holding that New Jersey's 2014 partial repeal ``authorize[d]'' sports gambling. Against the 1992 baseline in which sports betting was illegal across almost all states, repealing a prohibition both permits the activity and empowers private operators to conduct it. Neither the Third Circuit nor the United States could offer a coherent rule distinguishing permissible partial repeals from impermissible authorizations, and even under respondents' narrower reading the statute would still fail constitutional muster (\href{https://www.law.cornell.edu/supremecourt/text/16-476}{Cornell LII}).

Turning to the Tenth Amendment, Justice Alito grounded the anti-commandeering rule of \textit{New York v. United States}, 505 U.S. 144 (1992), and \textit{Printz v. United States}, 521 U.S. 898 (1997), in the Constitution's foundational break from the Articles of Confederation: Congress may regulate individuals directly, not state governments. He identified three core rationales: (1) dividing sovereignty serves as a structural safeguard of individual liberty; (2) prohibiting commandeering preserves political accountability so voters know which sovereign to credit or blame; and (3) it stops Congress from shifting regulatory and enforcement costs onto state budgets (\href{https://supreme.justia.com/cases/federal/us/584/16-476/}{Justia Opinion}).

Crucially, the majority rejected as ``empty'' the leagues' distinction between affirmative commands and negative prohibitions. Ordering a state legislature not to repeal an existing prohibition places it under direct congressional control no less than ordering it to pass a new law. Justice Alito distinguished four prior cases: \textit{South Carolina v. Baker}, 485 U.S. 505 (1988), and \textit{Reno v. Condon}, 528 U.S. 141 (2000), involved evenhanded regulations of activities in which both states and private actors participate; \textit{Hodel v. Virginia Surface Mining \& Reclamation Ass'n}, 452 U.S. 264 (1981), offered states a cooperative choice; and \textit{FERC v. Mississippi}, 456 U.S. 742 (1982), merely required state commissions to consider federal standards. He also rejected the Supremacy Clause preemption defense: all three forms of preemption---conflict, express, and field---operate through a federal law regulating private actors. Because \S~3702(1) neither grants private parties a federal right to gamble nor forbids private gambling directly, it is not preemption at all (\href{https://www.law.cornell.edu/supremecourt/text/16-476}{Cornell LII}).

In Part VI--B (for a six-Justice majority), Justice Alito held that no provision of PASPA was severable. Leaving intact the ban on state-run sports lotteries while permitting states to authorize private casinos would invert Congress's 1992 priorities. Leaving \S~3702(2) standing alone would produce a ``perverse policy'' making private sports gambling federally unlawful only where a state chooses to legalize it, contradicting traditional federal gambling laws and overwhelming a civil-injunction scheme that the Congressional Budget Office had scored at ``no cost'' to the federal government. Likewise, Congress would not have retained advertising bans on conduct legal under both federal and state law (\href{https://supreme.justia.com/cases/federal/us/584/16-476/}{Justia Opinion}).

\subsection{Concurring Opinion (Justice Thomas)}

Justice Thomas joined the Court's opinion in full while flagging two reservations. First, he expressed doubt that the Commerce Clause authorizes Congress to ban purely intrastate sports gambling even directly, citing \textit{License Tax Cases}, 72 U.S. (5 Wall.) 462 (1867), and his concurrence in \textit{United States v. Lopez}, 514 U.S. 549 (1995). Second, he urged the Court to revisit modern severability doctrine. Historically, early American courts simply declined to enforce unconstitutional provisions in the case before them without speculating about counterfactual legislative intent. Modern severability analysis, he warned, departs from textualism---because unexpressed hypothetical intent never passed bicameralism and presentment---and invites courts to pass judgment on statutory provisions that no litigant has Article III standing to challenge (\href{https://www.law.cornell.edu/supremecourt/text/16-476}{Cornell LII}).

\subsection{Concurrence in Part and Dissent in Part (Justice Breyer)}

Justice Breyer joined Parts I--V and VI--A, agreeing that \S~3702(1)'s bars on state authorization and licensing unconstitutionally commandeer state legislatures. He dissented from Part VI--B, however, concluding that \S~3702(2)---which directly prohibits private individuals from operating sports gambling schemes pursuant to state law---was a valid direct regulation of interstate commerce that could stand independently as a congressional fallback. Even if severing \S~3702(2) made New Jersey's victory ``mostly Pyrrhic,'' Justice Breyer reasoned that the constitutional flaw in \S~3702(1) lay in its means rather than Congress's permissible end (\href{https://supreme.justia.com/cases/federal/us/584/16-476/}{Justia Opinion}).

\subsection{Dissenting Opinion (Justice Ginsburg)}

Justice Ginsburg, joined in full by Justice Sotomayor and in part by Justice Breyer, focused on severability, assuming \textit{arguendo} that \S~3702(1)'s anti-authorization and anti-licensing clauses violated the anti-commandeering rule. She criticized the majority for using a ``wrecking ball'' and wielding ``an ax'' rather than a ``scalpel.'' Under \textit{Free Enterprise Fund v. PCAOB}, 561 U.S. 477 (2010), and \textit{New York}, courts should preserve valid portions of a statute unless it is evident Congress would have preferred no law at all. Leaving intact \S~3702(1)'s ban on state-operated schemes and \S~3702(2)'s ban on state-authorized private schemes would validly exercise Commerce Clause authority under \textit{Gonzales v. Raich}, 545 U.S. 1 (2005), without blurring political accountability, because voters would clearly attribute the prohibition to Congress (\href{https://www.law.cornell.edu/supremecourt/text/16-476}{Cornell LII}).

\section{Examples: Future Applications}

\begin{enumerate}
    \item \textbf{Same-Side (Unconstitutional Commandeering):} Congress enacts a statute stating that ``no State or municipality may prohibit its police officers from sharing immigration-status information with federal immigration agencies'' (mirroring challenges to 8 U.S.C. \S~1373). \textit{Why it falls on the same side:} Like PASPA \S~3702(1), this rule issues a direct negative order to state and local governments dictating what policies they may not adopt for their own officials, without regulating private conduct. Under \textit{Murphy}, forbidding a state from adopting a policy is commandeering just as much as forcing it to adopt one (\href{https://texaslawreview.org/making-federalism-great-again/}{Texas Law Review}; \href{https://hulr.org/spring-2022/murphy-v-ncaa}{HULR}).

    \item \textbf{Same-Side (Unconstitutional Commandeering):} Congress passes a statute providing that ``no State shall repeal or reduce its existing state-law criminal penalties for possession or distribution of cannabis.'' \textit{Why it falls on the same side:} Although Congress may directly criminalize cannabis under the federal Controlled Substances Act, it cannot lock state legislatures into maintaining state-law bans. Freezing existing state criminal codes directly controls state lawmaking and blurs voter accountability (\href{https://reason.com/volokh/2018/05/14/sports-gambling-decision-is-a-major-vict/}{Volokh Conspiracy}).

    \item \textbf{Opposite-Side (Constitutional Direct Regulation / Preemption):} Congress enacts a federal statute making it a federal felony for any private individual or company to operate a commercial sportsbook anywhere in the United States, enforced by federal prosecutors. \textit{Why it falls on the opposite side:} This statute directly regulates private actors rather than state legislatures. Although conflicting state laws yield under the Supremacy Clause, Congress neither commands state lawmakers nor shifts enforcement costs onto state governments (\href{https://supreme.justia.com/cases/federal/us/584/16-476/}{Justia Opinion}).

    \item \textbf{Opposite-Side (Constitutional Evenhanded Regulation):} Congress enacts a federal privacy law barring any entity---whether a private data broker or a state Department of Motor Vehicles---from selling consumers' biometric records without affirmative written consent. \textit{Why it falls on the opposite side:} Under \textit{Reno v. Condon} and \textit{South Carolina v. Baker}, reaffirmed in \textit{Murphy}, anti-commandeering does not apply when Congress evenhandedly regulates commercial activity in which both state agencies and private actors engage, rather than directing how states regulate their citizens (\href{https://www.law.cornell.edu/supremecourt/text/16-476}{Cornell LII}).

    \item \textbf{Fence-Sitter (Borderline / Hard Case):} Congress passes a statute granting private drone companies a federal statutory ``right to operate commercial delivery drones below 400 feet subject to FAA safety rules,'' coupled with a clause stating that ``no State or municipality may enact any zoning or privacy law restricting such flights.'' \textit{Why it is a fence-sitter:} On one hand, this resembles valid express preemption under \textit{Morales v. TWA} because Congress framed it as conferring an affirmative federal right on private actors subject only to federal constraints. On the other hand, if the FAA imposes virtually no substantive privacy or land-use standards, challengers could argue that Congress is impermissibly stripping states of core legislative power without enacting a genuine federal regulatory regime for private actors (\href{https://www.law.cornell.edu/supremecourt/text/16-476}{Cornell LII}).
\end{enumerate}

\section{Critique}

\textbf{Originalist and Textualist Perspectives:} Originalist and federalist scholars broadly praised \textit{Murphy} for closing a major loophole in Tenth Amendment doctrine. Writing on the \href{https://reason.com/volokh/2018/05/14/sports-gambling-decision-is-a-major-vict/}{Volokh Conspiracy}, Professor Ilya Somin commended Justice Alito for rejecting the formalistic divide between affirmative mandates and negative prohibitions, observing that without \textit{Murphy}, Congress could have eviscerated anti-commandeering simply by forbidding states from repealing prior laws. Somin also cited originalist work by Michael Rappaport arguing that the Founding-era meaning of ``State'' presupposed a sovereign polity whose legislative apparatus could not be commandeered by the federal government. At the same time, Justice Thomas's concurrence exposed an internal textualist tension: even while applying modern severability precedent to invalidate all of PASPA, textualists question whether judges have constitutional authority to speculate about ``hypothetical congressional intent'' instead of merely declining to enforce an unconstitutional provision in the case at hand (\href{https://www.law.cornell.edu/supremecourt/text/16-476}{Cornell LII}).

\textbf{Progressive and Doctrinal Perspectives:} Critics from progressive and nationalist traditions questioned both the textual basis of anti-commandeering and the sweep of the Court's severability ruling. In a \href{https://www.scotusblog.com/2017/08/symposium-time-abandon-anti-commandeering-dont-count-supreme-court/}{SCOTUSblog symposium}, Professor Steven D. Schwinn argued that anti-commandeering lacks explicit grounding in the Tenth Amendment, breeds unworkable line-drawing between preemption and commandeering, and should yield to political-process safeguards as in \textit{Garcia v. San Antonio Metropolitan Transit Authority}. Other commentators echoed Justice Ginsburg's critique that invalidating all of PASPA---including \S~3702(2)'s direct restriction on private actors---exceeded judicial restraint and functioned as a policy choice unleashing nationwide commercial sports gambling (\href{https://www.scotusblog.com/2018/05/opinion-analysis-justices-strike-down-federal-sports-gambling-law/}{SCOTUSblog}). Yet \textit{Murphy} also demonstrated federalism's shifting ideological valence: progressive cities and states immediately invoked \textit{Murphy}'s 7--2 anti-commandeering rule to defend ``sanctuary'' policies and state cannabis reforms against federal pressure (\href{https://texaslawreview.org/making-federalism-great-again/}{Texas Law Review}; \href{https://hulr.org/spring-2022/murphy-v-ncaa}{HULR}).

\textbf{Competing Constitutional Values:} At its core, \textit{Murphy} balances two competing structural ideals: \textit{state sovereignty, political accountability, and fiscal discipline} (preventing Congress from conscripting state legislatures to carry out federal objectives without bearing the political or financial cost) versus \textit{national regulatory flexibility and cooperative efficiency} (enabling Congress to set national policy floors without building an expansive federal enforcement apparatus).

\section{Key Quotations}

\begin{quote}
``A more direct affront to state sovereignty is not easy to imagine.''
--- Justice Samuel Alito, writing for the Court (\href{https://supreme.justia.com/cases/federal/us/584/16-476/}{Justia Opinion}).
\end{quote}

\begin{quote}
``The basic principle---that Congress cannot issue direct orders to state legislatures---applies in either event.''
--- Justice Samuel Alito, rejecting the affirmative-versus-negative command distinction (\href{https://www.law.cornell.edu/supremecourt/text/16-476}{Cornell LII}).
\end{quote}

\begin{quote}
``Congress can regulate sports gambling directly, but if it elects not to do so, each State is free to act on its own.''
--- Justice Samuel Alito, concluding the majority opinion (\href{https://www.law.cornell.edu/supremecourt/text/16-476}{Cornell LII}).
\end{quote}

\begin{quote}
``The Court wields an ax to cut down \S~3702 instead of using a scalpel to trim the statute.''
--- Justice Ruth Bader Ginsburg, dissenting on severability (\href{https://supreme.justia.com/cases/federal/us/584/16-476/}{Justia Opinion}).
\end{quote}

\end{document}
